\BeleskeID{04-s020}
% element: 04-s020:e01 — Pravci grupisanja celobrojnih i razlomljenih cifara od tačke
% element: 04-s020:e02 — Primer 1A.3C_16 sa prevođenjem i uklanjanjem vodećih/pratećih nula
% element: 04-s020:e03 — Primer 101100.1011011_2 sa dopunom grupa i rezultatom 2C.B6_16
% element: 04-s020:e04 — Oktalna analogija sa grupama po tri cifre
% element: 04-s020:d01

Tačka je granica od koje određujemo grupe u oba smera. Nepotpuna najviša celobrojna grupa dopunjava se levo, a nepotpuna poslednja razlomljena grupa desno. Tako sve postojeće cifre zadržavaju rastojanje od tačke i svoju težinu.

U drugom primeru razlomljena grupa $011$ dopunjuje se kao $0110$, a ne $0011$. Prvi zapis ima vrednost $6/16=3/8$, kao i početnih $0.011_2$; drugi bi imao vrednost $3/16$. Zbog toga dopuna poslednje grupe mora biti prateća nula. Pri uklanjanju nula smeju nestati samo vodeće nule celobrojnog i prateće nule razlomljenog dela.
